\chapter{Thesis Self-Assessment Report}

\section{General Profiles}
\begin{itemize}
    \item \textbf{Title of the Research:} A Hybrid Framework Using Explainable AI and Blockchain for Transparent Digital Document Verification
    \item \textbf{Name of Supervisor and Co-supervisor:} [Insert Your Supervisor's Name]
    \item \textbf{Synopsis/Abstract of the Research:} \textit{Refer to the translated Bengali/English abstract finalized in the previous conversation tier.}
\end{itemize}

\section{Mapping of Attained COs \& POs}


\begin{longtable}{|>{\centering\arraybackslash}m{1.2cm}|>{\centering\arraybackslash}m{1cm}|p{5.5cm}|>{\centering\arraybackslash}m{1cm}|>{\centering\arraybackslash}m{1cm}|>{\centering\arraybackslash}m{5cm}|}
\caption{Mapping of Attained COs \& POs} \label{tab:co_po_mapping} \\
\hline
\textbf{Term} & \textbf{COs} & \textbf{CO Statements} & \textbf{POs} & \textbf{Level} & \textbf{Relevant Chapter / Section / Article} \\ \hline
\endfirsthead
\caption[]{Mapping of Attained COs \& POs (Continued)} \\
\hline
\textbf{Term} & \textbf{COs} & \textbf{CO Statements} & \textbf{POs} & \textbf{Level} & \textbf{Relevant Chapter / Section / Article} \\ \hline
\endhead

\textbf{Spring} & \textbf{CO1} & Identify a real-life problem / contemporary issue that can be translated to an engineering solution focusing on potential research innovation. & \textbf{PO12} & \textbf{H} & Chapter 1: Section 1.1, 1.2, 1.3 \\ \hline
\textbf{Spring} & \textbf{CO2} & Identify outcomes and functional requirements of the proposed solution considering software and/or hardware specifications and standards. & \textbf{PO2} & \textbf{H} & Chapter 3: Section 3.1, 3.2, 3.5 \\ \hline
\textbf{Spring} & \textbf{CO4} & Investigate and analyze engineering/computing systems with given specifications and requirements. & \textbf{PO4} & \textbf{H} & Chapter 4: Section 4.2, 4.4; Chapter 5: Section 5.2, 5.3 \\ \hline
\textbf{Spring} & \textbf{CO5} & Design and build the proposed solution of the problem considering the requirements and evaluate the system. & \textbf{PO3} & \textbf{H} & Chapter 4: Section 4.3, 4.5; Chapter 5: Section 5.4 \\ \hline
\textbf{Fall}   & \textbf{CO1} & Use modern analysis and design tools in the process of designing and validating of a system/subsystem. & \textbf{PO5} & \textbf{M} & Chapter 4: Section 4.1 (Table 4.1); Appendices B, C, G \\ \hline
\textbf{Fall}   & \textbf{CO2} & Assess social impacts and responsibilities of the engineering design project. & \textbf{PO6} & \textbf{M} & Chapter 1: Section 1.5; Annex A: Section 1 (Social Constraints) \\ \hline
\textbf{Fall}   & \textbf{CO3} & Evaluate the ethical and professional implications of the proposed solution. & \textbf{PO8} & \textbf{H} & Chapter 3: Section 3.4; Annex A: Section 1 (Ethical Considerations) \\ \hline
\textbf{Fall}   & \textbf{CO4} & Identify and validate the impact of environmental considerations and the sustainability of a system/subsystem. & \textbf{PO7} & \textbf{L} & Chapter 4: Section 4.3.2; Annex A: Section 1 (Environmental) \\ \hline
\end{longtable}

\section{Assessment Questions and Answers}
\begin{enumerate}[label=\roman*.]
    \item \textbf{What real-world problem did you address, what research gap did you identify, and how does your solution contribute to future innovation or advancement in the field?} \\
    \textbf{Answer:} I addressed the worldwide proliferation of academic credential fraud and the rapid emergence of high-fidelity, synthetic (AI-generated) certificate forgeries. The core research gap identified was that traditional blockchain implementations (such as \textit{BCert} or \textit{BeCertify}) act merely as passive storage vaults—they lack ``pre-validation'' mechanisms, meaning a fraudulent certificate uploaded to the network becomes permanently legitimized on-chain. My solution introduces a Python-driven Explainable AI (XAI) verification tier that audits linguistic features (perplexity/burstiness) and layout metadata before transaction anchoring occurs, transforming opaque ``black-box'' decisions into trustworthy ``glass-box'' validations.

    \item \textbf{How did you analyze the problem requirements, identify constraints, and determine the functional specifications of your proposed solution?} \\
    \textbf{Answer:} I systematically categorized requirements across three distinct domains: data extraction, computational analysis, and decentralized tracking. Constraints centered on data safety (avoiding plain-text sensitive info on-chain), storage costs (gas consumption optimization), and explanation clarity. The functional requirements were determined as: complete parsing of 15+ file types, SHA-256 fingerprint matching, multi-dimensional semantic vector indexing via Neon PostgreSQL, and real-time smart contract state queries.

    \item \textbf{How did you plan, schedule, allocate resources, and manage risks throughout the project to ensure successful completion?} \\
    \textbf{Answer:} The project was planned using a three-phase software lifecycle. Phase 1 built a modular proof-of-concept for core hashing and Solidity mapping; Phase 2 advanced the Minimum Viable Product (MVP) via off-chain Neon integration and relational user routing; Phase 3 automated child-process pipelines. Principal risks involved layout extraction errors and transaction execution bottlenecks. These were managed by utilizing parallel scripts that trigger the vector search and AI detector simultaneously to curb computational blockages.

    \item \textbf{What investigation and analysis methods did you use to validate your assumptions, and how did the results influence your solution?} \\
    \textbf{Answer:} I constructed an empirical experimental test dataset comprising 200 documents, meticulously distributed between genuine files, structural layout adjustments, synthetic text generators (GPT-4/Claude), and rephrased plagiarism blocks. The empirical findings confirmed that the model achieved a 92.5\% classification rate for synthetic content and isolated 48 out of 50 layout anomalies. These results verified my core hypothesis: statistical evaluation metrics (low perplexity/zero sentence variation) remain highly effective indicators of synthetic forgery.

    \item \textbf{How did you design, implement, and evaluate your solution to ensure that it satisfies the identified requirements and constraints?} \\
    \textbf{Answer:} The solution was designed around a split-pipeline model. Relational profiles and high-dimensional semantic embeddings (768-dimensional text vectors generated via BERT) are stored in the serverless cloud (Neon DB), while the core validation proofs (SHA-256 signatures) are routed via ethers.js onto an Ethereum smart contract ledger. Evaluation benchmarks demonstrated a transaction confirmation rate of 12.4 seconds, with retrieval queries returning values in under 2 seconds, effectively confirming the setup's operational scalability.

    \item \textbf{Which modern engineering or computing tools did you use, why were they selected, and how did they support the development and validation of your project?} \\
    \textbf{Answer:} I used a specialized technical ecosystem to drive validation:
    \begin{itemize}
        \item \textit{Node.js \& Express:} Chosen for high-throughput, non-blocking I/O traffic management.
        \item \textit{Solidity \& Hardhat:} Used to program and compile robust local contract states.
        \item \textit{Python (PyTorch/Transformers):} Critical for computing language metrics.
        \item \textit{Neon PostgreSQL \& pgvector:} Selected to facilitate rapid mathematical similarity evaluations (cosine distances) under 500ms, surpassing standard SQL queries.
    \end{itemize}

    \item \textbf{What are the societal impacts of your project, and how does your solution address the needs and responsibilities of stakeholders?} \\
    \textbf{Answer:} By establishing a reliable framework, this project systematically mitigates credential counterfeiting, protecting the hard work of legitimate students and minimizing financial overhead for hiring institutions. It fulfills an important societal responsibility by shifting background screenings away from time-consuming, expensive central dependencies toward direct, human-interpretable, decentralized validation architectures.

    \item \textbf{What ethical and professional issues are associated with your project, and how did you address them during the design and implementation process?} \\
    \textbf{Answer:} Opaque automation risks violating an individual's ethical right to an explanation when their certificate is flagged. My system embeds explainability report models directly into the pipeline, giving users interactive evidence maps of flagged segments. Professionally, I adhered to data privacy restrictions by ensuring all persistent relational entities are contained securely on cloud databases, using only anonymized hashes on public ledgers.

    \item \textbf{How does your project consider environmental sustainability, and what measures were taken to minimize negative environmental impacts?} \\
    \textbf{Answer:} Extensive block storage commands represent massive electricity costs across mining structures. To align with environmental goals, my implementation stores only a 32-byte hash footprint on-chain. Large parsed files are delegated to decentralized P2P networks (IPFS via Pinata), dramatically lowering resource extraction footprints on the Ethereum Virtual Machine.

    \item \textbf{What was your role in the team, and how did collaboration among team members contribute to achieving the project objectives?} \\
    \textbf{Answer:} \textit{[Customize this according to your specific work role, e.g.,]} As the lead systems developer, I spearheaded the technical engineering of the Express routing orchestrator and directed the backend integration with the serverless Neon PostgreSQL cluster. My coordination with my co-researchers ensured that the Python-driven NLP metrics mapped precisely onto the Solidity smart contract schemas, allowing our distributed microservices to execute seamlessly.

    \item \textbf{How would you communicate your project's objectives, methodology, results, and significance to both technical and non-technical audiences?} \\
    \textbf{Answer:} For technical audiences, I rely on rigorous code metrics, schema maps, gas optimization limits, and perplexity curves. For non-technical audiences, I use a clear structural analogy: the system serves as a digital airport checkpoint where a smart automated guard inspects the file's text signatures, an AI detective explains the red flags, and an immutable ledger permanently saves the receipt so that it can never be altered or forged.
\end{enumerate}